%% kapa: exam=1 round=23 year=2012 period=1 form=A subject=economics subject_ko=경제원론 seq=21 number=61
%% source: Q-Net 감정평가사 기출문제 (cst00309) artlSeq=5052578 file=raw/1008/5052578/2012년 제23회 감정평가사 제1차시험 기출문제(1교시 A형).pdf
%% conversion: kordoc 4.24.51 -> tools/gen1.py (auto); review=proofread-sonnet; flags=none
\begin{question}{61}
개방경제의 국민소득계정에 관한 설명으로 옳은 것을 모두 고른 것은?
\begin{framedtext}
ㄱ. 국민소득이 소비, 투자, 정부지출의 합보다 큰 경우에 순수출은 반드시 양(+)이 된다. \par
ㄴ. 민간 투자가 민간 저축보다 더 큰 경우에 순수출은 반드시 양(+)이 된다. \par
ㄷ. 정부 세금 수입이 지출보다 더 큰 경우에 순수출은 반드시 양(+)이 된다.
\end{framedtext}
\begin{choices}
  \item ㄱ
  \item ㄴ
  \item ㄷ
  \item ㄱ, ㄴ
  \item ㄴ, ㄷ
\end{choices}
\end{question}
\answer{1}
